\documentclass[11pt]{article}
\usepackage[T1]{fontenc}\usepackage{lmodern}
\usepackage[a4paper,margin=24mm]{geometry}
\usepackage{amsmath,amssymb,amsthm,mathrsfs,mathtools,booktabs,longtable,array,microtype}
\usepackage[hidelinks]{hyperref}\usepackage{xurl}
\newtheorem{theorem}{Theorem}[section]
\newtheorem{proposition}[theorem]{Proposition}
\newtheorem{lemma}[theorem]{Lemma}
\newtheorem{corollary}[theorem]{Corollary}
\theoremstyle{definition}\newtheorem{definition}[theorem]{Definition}
\theoremstyle{remark}\newtheorem{remark}[theorem]{Remark}
\newcommand{\Z}{\mathbb Z}\newcommand{\Q}{\mathbb Q}\newcommand{\R}{\mathbb R}
\newcommand{\N}{\mathbb N}\newcommand{\ES}{\operatorname{ES}}
\newcommand{\Div}{\operatorname{Div}}\newcommand{\supp}{\operatorname{supp}}
\newcommand{\lcm}{\operatorname{lcm}}\newcommand{\Tr}{\operatorname{Tr}}
\DeclareMathOperator{\Res}{Res}
\title{Positive witnesses and a constructive counterfactual system\\for the Erd\H{o}s--Straus equation}
\author{The Clankers research workbench\\\small Assistant-developed continuation; source-specific credit retained}
\date{12 September 2026}
\begin{document}\maketitle
\begin{abstract}
The same full finite shell atlas supports two exact research directions. Every
canonical divisor state gives a positive rational decomposition of $4/p$;
the two nondistinguished denominators have the same reduced denominator $D$.
We compute $D$ exactly, identify failure at $p$ with an explicitly invertible
finite diagonal operator, and carry both successful and unsuccessful states
through marked Cayley, factor, Leech and octonionic coordinates. Integer
collision-energy bounds and an exact largest non-subgroup-divisor calculation
supply constructive sufficient conditions for $D=1$, including a strict
improvement on the preceding angular bound. No counterexample is found and
no universal positivity theorem is claimed. The all-shell criterion itself,
the E/M parametrizations, and the standard inequalities are inherited;
the present contribution is the explicit two-sided continuation and its
certificates, without a historical-priority assertion.
\end{abstract}

\section{Scope, source lineage, and the Giuga model}
The target is $4/n=1/x+1/y+1/z$ with positive integer denominators.
A counterexample to a restricted ray family, an angular chamber, or the
cutoff $R<p$ is not a counterexample to this equation. We use the entire
first-denominator range forced by a smallest denominator. The foundation's
\texttt{all\_shell\_no\_hit.tex}, pinned at
\texttt{9a3ddfe8b38a424f1ff0426201a1f4f6f9d2772f}, already supplies the
finite exponent packet and complete ordered return \cite{foundation}.
The antecedents are Yamamoto's Lemma 2, equation (4), p.~38;
Monks--Velingker's reduced-fraction criterion, Theorem 2.2, pp.~4--6;
and Elsholtz--Tao's Type I/II maps, \S2, Propositions 2.2 and 2.6
\cite{yamamoto,mv,et}. All formulas required here are rederived below.

The uploaded \texttt{verify\_giuga\_blind\_plane\_cd.py} separates a Giuga
integer from a counterexample to Giuga's conjecture. Its blocks A3--A6
retain the additional Carmichael conditions. For a composite candidate $n$
and each prime $q\mid n$, the simultaneous local conditions are
\[
 q\mid n/q-1,\qquad q-1\mid n-1
 \quad\Longleftrightarrow\quad q^2(q-1)\mid n-q.
\]
The equivalence follows from $q^2\mid n-q$ and
$n-q\equiv n-1\pmod{q-1}$, since $\gcd(q^2,q-1)=1$.
The first condition forces squarefreeness. The counterexample must be both
Giuga and Carmichael, not merely one of them; see Kellner's Theorem 3.5
\cite{kellner}. We import this \emph{method of retaining simultaneous
conditions}, not a conjectural map between Giuga and ES counterexamples.
The file's A7 uses floating complex exponentials, while A9b checks an exact
CRT bijection; neither is represented here as a full formal proof.

For exact phase conventions, the pair
$(e^{2\pi iA/q},e^{2\pi iB/(q-1)})$ has relative exponent
\[
 e=(q-1)A-qB\pmod{q(q-1)}.
\]
Its inverse is $A=-e\pmod q$, $B=-e\pmod{q-1}$. Thus this restricted
relative phase is faithful, although a generic equal-radius Hopf quotient
forgets the common phase. This fixes the signs explicitly without relying
on numerical equality of roots of unity. The source's generic annihilator
and Cayley--Dickson experiments are not needed for the ES proof below.

\section{The full candidate atlas}
Fix a prime $p=4m+1$, $m\ge1$, and set
\[
 \mathcal S_p=\{m+1,\ldots,3m\},\qquad R_a=4a-p.
\]
Then $3\le R_a\le2p-3$, $R_a\equiv3\pmod4$, and
$\gcd(a,R_a)=\gcd(p,R_a)=1$. These follow from $0<a<p$ and primality.
Write $a=\prod_q q^{e_q}$ and retain the whole box
\[
 \mathcal B_a=\prod_{q\mid a}\{0,\ldots,2e_q\},\qquad
 u(f)=\prod_{q\mid a}q^{f_q}.
\]
Unique factorization makes this a bijection to $\Div(a^2)$, with inverse
$f_q=v_q(u)$. No generated subgroup replaces this box.

A canonical candidate is $\alpha=(p,a,u,c)$ with $a\in\mathcal S_p$,
$u\mid a^2$, and $c\in\{E,M\}$. Define
\begin{align}
 Q_E(\alpha)&=\left(a,\frac{pa+a^2/u}{R_a},\frac{pa+p^2u}{R_a}\right),\label{eq:E}\\
 Q_M(\alpha)&=\left(a,\frac{p(a+a^2/u)}{R_a},\frac{p(a+u)}{R_a}\right).\label{eq:M}
\end{align}
The initial position is distinguished and the other two positions remain
ordered. An additional exterior swap bit is used when all raw orders are
required. A middle swap is already $u\mapsto a^2/u$; adding another middle
swap bit would duplicate records.

\begin{theorem}[Positive rational atlas and complete integral test]\label{thm:atlas}
Every candidate in \eqref{eq:E}--\eqref{eq:M} is a positive rational
solution of $4/p=1/a+1/Y+1/Z$. With
\[
 g=\gcd(a,u),\quad r=u/g,\quad s=a/g,\quad h=g^2/u,
\]
one has $h,r,s\in\Z_{>0}$, $\gcd(r,s)=1$, $a=hrs$, and $u=hr^2$.
The rational quotients are
\[
 \kappa=(pr+s)/R_a\quad(E),\qquad
 \lambda=(r+s)/R_a\quad(M).
\]
The ordered maps are $(hrs,hs\kappa,phr\kappa)$ and
$(hrs,phs\lambda,phr\lambda)$. Moreover,
\[
 \ES(p)\quad\Longleftrightarrow\quad
 \text{at least one canonical candidate has all coordinates integral}.
\]
There are at most $(p-1)^2$ canonical candidates.
\end{theorem}
\begin{proof}
At a prime with $v_q(a)=e$ and $v_q(u)=f$, the exponents of
$r,s,h$ are $(f-e)_+,(e-f)_+$ and $e-|f-e|$. They are nonnegative
and prove all normalization assertions. Direct substitution gives
\[
 (R_aY-pa)(R_aZ-pa)=p^2a^2
\]
in both channels. Expanding yields
$R_aYZ=pa(Y+Z)$, the desired identity; every displayed denominator is
positive. The quotient formulas follow by substitution.

For completeness take an actual integral solution and move its smallest
denominator to the first position, retaining that permutation.
Positivity gives $p/4<a\le3p/4$, hence $a\in\mathcal S_p$.
Put $d_Y=R_aY-pa$, $d_Z=R_aZ-pa$. Then $d_Y,d_Z>0$ and
$d_Yd_Z=p^2a^2$. Thus $d_Y=p^iv$ for $i\in\{0,1,2\}$,
$v\mid a^2$. If $i=0$, the inverse is $E,u=a^2/d_Y$ without a swap;
if $i=1$, it is $M,u=pa^2/d_Y$; if $i=2$, it is
$E,u=d_Y/p^2$ with the exterior swap. The congruences
$d_Y\equiv d_Z\equiv-pa\pmod{R_a}$ give exactly the required channel
gates. For instance $i=0$ gives $R_a\mid4u+1$ after multiplication by
units $u,a^{-1}$ and use of $p\equiv4a$; $i=1$ gives
$R_a\mid u+a$. Conversely an integral candidate is already a witness.
The explicit inverse also preserves the full divisor exponent vector.

Pair divisors of $a^2$ around its square root $a$ to get
$\tau(a^2)\le2a-1$. Hence
\[
 |\mathscr A_p|=2\sum_{a=m+1}^{3m}\tau(a^2)
 \le2\sum_{a=m+1}^{3m}(2a-1)=16m^2=(p-1)^2.
\]
The bound counts canonical candidates, not their additional exterior swaps.
\end{proof}

\section{The exact return denominator}
\begin{theorem}[Common denominator and defect]\label{thm:denom}
The two nondistinguished coordinates of each candidate have the same
reduced rational denominator
\[
 D_E=\frac{R_a}{\gcd(R_a,4u+1)},\qquad
 D_M=\frac{R_a}{\gcd(R_a,u+a)}.\tag{D}
\]
This $D$ is the least positive integer clearing the entire triple.
It is odd, divides $R_a$, and is prime to $p$. With
\[
 X=DQ_c(\alpha),\qquad n=pD,\qquad
 \delta_\alpha=(D-1)/2,
\]
one has $X\in\Z_{>0}^3$, $\sum 1/X_i=4/n$, and
\[
 F_p(X):=4X_0X_1X_2-p\sum_{i<j}X_iX_j
 =2p \delta_\alpha\sum_{i<j}X_iX_j.\tag{F}
\]
Success at the original prime is exactly $D=1$.
If $D>1$, then $3p\le n\le p(2p-3)$ and
$1\le \delta_\alpha\le p-2$.
\end{theorem}
\begin{proof}
For $E$, multiply the numerator of $Y$ by the unit $u$ modulo $R_a$:
\[
 u(pa+a^2/u)\equiv a^2(4u+1),\qquad
 pa+p^2u\equiv4a^2(4u+1)\pmod{R_a}.
\]
Every displayed multiplier $u,a,4$ is a unit. Thus both numerator gcds
with $R_a$ equal $\gcd(R_a,4u+1)$. For $M$, the numerator identities
\[
 u\,p(a+a^2/u)=pa(u+a),\qquad p(a+u)=p(u+a)
\]
give the same conclusion with $u+a$. Reducing the fractions proves (D).
As $a$ is integral, $D$ is the least clearing factor of all three entries.
The bounds and coprimality follow from the shell units. Finally
$F_{pD}(X)=0$ by scaling the reciprocal identity; subtracting it from
$F_p(X)$ gives (F).
\end{proof}

\begin{proposition}[Complete candidate inverse and collision stratum]
In the canonical order,
\[
 u=\frac{a^2}{R_aY-pa}\ (E),\qquad
 u=\frac{pa^2}{R_aY-pa}\ (M).
\]
An exterior triple has three distinct entries. A middle triple has a
repeat precisely when $u=a$, in which case $Y=Z$ and $D=R_a>1$.
The middle complement preserves $D$ and reverses the last two entries.
\end{proposition}
\begin{proof}
The residual products in the two channels are respectively
$(a^2/u,p^2u)$ and $(pa^2/u,pu)$, proving the inverse.
In the exterior channel $v_p(Y)=0$ and $v_p(Z)\ge1$;
$a=Y$ would give $(4u-1)a=2pu$. Since $\gcd(u,4u-1)=1$,
$4u-1$ divides $2p$, impossible modulo four for its odd divisor.
The equality $a=Z$ would imply $p\mid4a^2$, also impossible.
In the middle channel both $Y,Z$ have positive $p$-adic valuation,
while $a$ does not. The equality $Y=Z$ is $a^2/u=u$, hence $u=a$.
Then $\gcd(R_a,2a)=1$ gives $D=R_a$.
Complementation follows from the displayed formulas.
\end{proof}

\begin{proposition}[A same-shell map cannot improve this defect]\label{prop:transport}
If $u,u'\mid a^2$ and $u'\equiv pu\pmod{R_a}$, the exterior
candidate at $u$ and the middle candidate at $u'$ have exactly the same $D$,
whether or not they are integral.
\end{proposition}
\begin{proof}
$u'+a\equiv a(4u+1)\pmod{R_a}$, and $a$ is a unit. Apply (D).
\end{proof}
This extends the inherited same-shell congruence to unsuccessful rational
states. It does not assert that a suitable $u'$ always exists, nor does
failure of this particular improvement mechanism rule out other maps.

\section{The counterfactual object and its finite inverse certificate}
\begin{definition}
The \emph{defect atlas} at $p$ consists of all shells, their complete prime
factorizations and exponent boxes, both canonical channel maps, each $D$,
and each local valuation deficit
\[
 v_\ell(D)=\max\{v_\ell(R_a)-v_\ell(g_c(u)),0\},
 \quad g_E(u)=4u+1,\quad g_M(u)=u+a.
\]
A negative certificate lists every canonical candidate and, for each one,
a prime $\ell$ where the displayed deficit is positive. It has the
explicit scope \emph{full} or \emph{restricted}.
\end{definition}

Let $\mathcal A_p=\Q^{\mathscr A_p}$ be the finite algebra of functions on
the canonical candidate set, retaining its basis labels. Let $H_p$ be
multiplication by $\delta_\alpha=(D_\alpha-1)/2$.
\begin{theorem}[Two-sided certificate theorem]\label{thm:dual}
The following are equivalent:
\begin{enumerate}
\item $p$ has no positive integral ES decomposition;
\item $D_\alpha\ge3$ for every canonical candidate;
\item $H_p$ is invertible;
\item there is a unique function $b$ with $\delta b=1$, namely
$b(\alpha)=2/(D_\alpha-1)$;
\item there is a full negative certificate with the stated local deficits.
\end{enumerate}
When these conditions hold, $0<b(\alpha)\le1$. Otherwise any zero coordinate
of $H_p$ reconstructs a positive integral witness. Both alternatives are
recognizable by terminating exact algorithms.
\end{theorem}
\begin{proof}
The atlas theorem and the common denominator theorem identify existence
with some $D=1$. Every other $D$ is odd and at least three. A diagonal
operator on a finite coordinate algebra is invertible exactly when all
entries are nonzero, with the displayed unique inverse. A nontrivial
positive integer $D$ has a prime divisor, whose valuation has the stated
positive deficit. Conversely any such deficit proves $D>1$.
Enumerating the finite shell interval and all bounded exponent choices
constructs one of the alternatives and verifies it. Primality is checked
by trial division through $\lfloor\sqrt p\rfloor$ in the supplied verifier.
\end{proof}

There is a universal polynomial version, not an uncomputed hypothesis.
Set
\[
 \Pi_p(z)=\prod_{j=1}^{p-2}(1-z/j),\qquad
 Q_p(z)=\frac{1-\Pi_p(z)}z\in\Q[z].
\]
All diagonal entries of $H_p$ belong to $\{0,\ldots,p-2\}$, so
\[
 \Pi_p(H_p)=\text{the exact success projector},\qquad
 H_pQ_p(H_p)=I-\Pi_p(H_p).
\]
Thus a false prime has $H_pQ_p(H_p)=I$. The success dimension is the number
of canonical passing candidates, not a count of globally sorted triples.
This interpolation is a standard finite-algebra fact; the particular
operator and its arithmetic labels are the construction here.

\paragraph{A genuine restricted prototype, not a counterexample.}
At $p=241,a=64,R=15$ all divisors are $u=2^e$, $0\le e\le12$.
For either channel the $D$ histogram is
\[
 D=3:4\text{ entries},\quad D=5:6\text{ entries},\quad D=15:3\text{ entries}.
\]
The negative operator is invertible on these 26 slots. Yet the \emph{full}
atlas at $241$ has positive witnesses in other shells. Even
$\gcd\{D_\alpha\}=1$ on this rejected shell: taking a gcd across different
candidates would give a false conclusion. The second prototype is
$p=37,a=18,R=35$, whose integral exponent lift can lie outside the original
box. Both prototypes are already motivated in the foundation and are
replayed here with their full rational defect data.

The same example has a complete defect-preserving transport graph. Its
vertices are the 26 labelled $E/M$ candidates. Join complementary middle
states, and join $E(u)$ to $M(u')$ precisely when
$u'\equiv pu\pmod R$. Every edge preserves $D$ by
Proposition~\ref{prop:transport}. Exhaustive adjacency gives exactly
three connected components: eight vertices at $D=3$, twelve at $D=5$,
and six at $D=15$. Thus combining these particular moves cannot turn
this rejected shell into a passing shell. The graph does not restrict
other cross-shell moves and is not a global counterexample.

\paragraph{Consequences for a least global counterexample.}
A least counterexample $n$ must be prime: multiply a solution at any prime
divisor by $n/p$. The prime $2$ is solved by $(1,2,2)$. For $p\equiv3\pmod4$
the triple $((p+1)/2,(p+1)/2,p(p+1)/4)$ works. For $p\equiv2\pmod3$ use
$(p,(p+1)/3,p(p+1)/3)$. A remaining prime is $1\pmod{12}$; if it is
$13\pmod{24}$ then $a_3=(p+3)/4$ is even and $u=2$ passes the residual-three
exterior gate. Therefore a least false prime is $1\pmod{24}$.
Its first two shells obey
\[
 q\mid (p+3)/4\Longrightarrow q\equiv1\pmod3,
\]
\[
 q\mid (p+7)/4\Longrightarrow q\equiv1,2,4\pmod7.
\]
The first assertion is the exact residual-three criterion. For the second,
$a_7$ is even and the signed powers $2^{-1},1,2$ fill the subgroup
$\{1,2,4\}$ modulo seven. Any factor outside it makes the middle target
$-1$ attainable using its one actual occurrence and one of those three
powers. If every factor is inside, both original targets are outside.
These are specializations of the foundation's first-two-shell theorems,
not new source attributions \cite{foundation}.

The condition $p+1$ having a prime divisor $q\equiv3\pmod4$ also gives a
witness: $R=q<p$, $r=s=1$, $h=(p+q)/4$, and $q\mid p+1$.
Thus a false prime has no such factor. Every further necessary condition
proved below applies to every shell of the same false prime.

\section{A sharpened positive collision criterion}
Fix a shell and $G=(\Z/R\Z)^\times$. Use the actual signed exponent box
\[
 \mathcal B^\pm_a=\prod_{q\mid a}\{-e_q,\ldots,e_q\}.
\]
Let $\mu(g)$ count its full fibres under
$\beta\mapsto\prod q^{\beta_q}\pmod R$, and put
\[
 N=\prod_q(2e_q+1),\quad E=\sum_g\mu(g)^2,
 \quad\mathcal T=\{-1,-p,-p^{-1}\}\subset G.
\]
The set braces remove any coincident targets. Inversion of the box gives
$\mu(g)=\mu(g^{-1})$. The channel targets on $u/a$ are $-p^{-1}$ and $-1$;
therefore avoiding both channels is exactly avoiding $\mathcal T$.

\begin{theorem}[Integer collision-energy forcing]\label{thm:energy}
Let $M=|G|-|\mathcal T|$. If $M>0$ write $N=Mq+r$ with $0\le r<M$ and set
\[
 L(N,M)=Mq^2+r(2q+1).
\]
If both channels are empty, then $E\ge L(N,M)$.
In particular $E<L(N,M)$ constructs a witness. If $M=0$, a witness exists.
The exact collision energy is
\[
 E=\sum_{\delta_q=-2e_q}^{2e_q}
 \mathbf1_{\prod q^{\delta_q}=1\ (R)}
 \prod_q(2e_q+1-|\delta_q|).\tag{C}
\]
\end{theorem}
\begin{proof}
A miss leaves at most $M$ bins containing $N$ occurrences. If two bin
counts differ by at least two, moving one from the larger to the smaller
strictly decreases the sum of squares. The minimum therefore has $r$
bins of size $q+1$ and $M-r$ bins of size $q$, giving $L$.
For (C), $E$ counts ordered pairs of exponent vectors with equal image.
At a fixed difference $\delta_q$, exactly $2e_q+1-|\delta_q|$ pairs
occur in that coordinate; multiply and impose the one simultaneous
residue equation. To extract a hit, retain an exponent vector in a target
fibre. At $-1$ use $M$. At $-p^{-1}$ use $E$. At $-p$ negate the whole
exponent vector first. In each case $u=\prod q^{e_q+\beta_q}$ remains
inside the original divisor box and the canonical denominator formula
returns the ordered witness.
\end{proof}

\paragraph{A stronger exact symmetric minimum.}
If $1\in\mathcal T$, the zero exponent vector is already a hit.
Otherwise let $s$ be the number of allowed nonidentity involutions and
$r$ the number of allowed unordered inverse pairs. The identity count is
odd, each other involution count is even, and each inverse pair has equal
counts. A miss consequently satisfies
\[
 E\ge L_{\rm sym}:=
 \min_{x_0+\cdots+x_s+y_1+\cdots+y_r=(N-1)/2}
 \left((1+2x_0)^2+4\sum_{i=1}^sx_i^2+2\sum_{j=1}^ry_j^2\right),
\]
with nonnegative integer variables. This finite minimum has an exact greedy
algorithm: start from one and take $(N-1)/2$ least available marginal costs
from the increasing sequences $8,16,24,\ldots$ for the identity,
$4,12,20,\ldots$ for each other involution, and $2,6,10,\ldots$ for each
inverse pair. Each selected prefix is feasible, and an exchange of an
unnecessarily expensive selected marginal with a cheaper available one
proves optimality. These constraints follow from the involution
$\beta\mapsto-\beta$, whose sole fixed vector is zero.

\paragraph{Strict improvement at an inherited stress prime.}
For $p=1201,a=306,R=23$, the factors are $2\cdot3^2\cdot17$.
Here $N=45$, $\mathcal T=\{9,18,22\}$, $M=19$, and the exact energy is
$107$. Ordinary Cauchy--Schwarz gives only
$19\cdot107=2033>2025$ and does not force a hit.
The integer minimum is $12\cdot2^2+7\cdot3^2=111$;
$L_{\rm sym}=111$ as well. Since $107<111$, the sharper theorem forces
an actual target. The retained inverse returns $u=17$, channel $E$,
$(h,r,s,\kappa)=(17,1,18,53)$ and
$(306,16218,1082101)$. This is not new finite coverage of that prime;
it is a strictly stronger exact sufficient criterion.

A hypothetical false prime must satisfy $E_a\ge L_{{\rm sym},a}$ in
\emph{every} full shell. Equivalently the nontrivial Fourier mass of each
multiplicity packet must be sufficiently large, since the usual finite
Parseval identity gives
$\sum_{\chi\ne1}|\widehat\mu(\chi)|^2=|G|E-N^2$.
No Fourier sign, density estimate, or subgroup-support approximation
replaces the integer target fibre.

\section{Exact maximal-cofactor angular forcing}
Let $H\le(\Z/R\Z)^\times$ have index two with $-1\notin H$.
Let $A\mid a$, all of whose prime factors lie in $H$, and retain the
\emph{actual} finite anchor pairs
\[
 (n,d):\quad nd\mid A,\quad\gcd(n,d)=1.
\]
Suppose their ratios cover $H$, an explicitly checkable condition on $A$.
For $v\in H$, define the attained cost
\[
 c_A(v)=\min_{n/d=v\ (R)}d/n,
 \qquad B_A=\max_{v\in H}c_A(v).
\]
This is the preceding angular construction's exact finite data, not a
universal assumption on the unknown ES prime.

\begin{theorem}[Exact largest available outside divisor]\label{thm:cofactor}
Write $a=Ab$ and assume $\gcd(a,R)=1$. The largest divisor of $b$ outside $H$
is
\[
 T_H(b)=
 \begin{cases}
 b,&b\notin H,\\
 b/q_*,&b\in H\text{ and has a prime divisor outside }H,\\
 0,&\text{every prime divisor of }b\text{ lies in }H,
 \end{cases}
\]
where $q_*$ is the least actual outside prime factor.
If $T_H(b)>8B_A$, there is a middle state with $0<r/s<1/8$.
More exactly, the anchor construction has such a state if and only if
some actual outside divisor $t\mid b$ satisfies
$t>8c_A(-t^{-1})$.
\end{theorem}
\begin{proof}
If $b$ is outside, $b$ itself is the largest divisor. If $b\in H$ and $t$
is outside, the complement $b/t$ is outside, hence has an outside prime
factor and is at least $q_*$. The bound $t\le b/q_*$ is attained by that
quotient. If there is no outside factor no divisor can be outside.

For an outside divisor $t$, choose a minimizing pair $(n,d)$ for
$-t^{-1}\in H$. Then $R\mid d+tn$ and $dtn\mid a$. Set
\[
 g=\gcd(d,tn),\quad r=d/g,\quad s=tn/g,\quad
 h=\frac{ag^2}{dtn},\quad u=\frac{ad}{tn},\quad
 \lambda=\frac{d+tn}{gR}.
\]
Every coordinate is integral; $g$ is a unit modulo $R$;
$a=hrs$, $u=hr^2$, and $\gcd(r,s)=1$. Thus this is a middle state,
and $r/s=c_A(-t^{-1})/t$. This proves both assertions. For a marked target
$(h,r,s)$ the complete inverse fibre consists of the anchor pairs for which
$t=ds/(nr)$ is a positive divisor of $b$ outside $H$ and satisfies the
stated congruence and inequality. Coprimality forces $r\mid d$ and then
$g=d/r$. Thus all inverse data are explicit; forgetting $(n,d,t)$ can
create only the stated finite fibres.
\end{proof}
The lower Koide boundary exceeds $1/8$: the rational bounds
$\sqrt3>1732/1000$, $\sqrt2<1415/1000$, give
$317/2415>1/8$. The radical bounds follow by squaring positive rationals.
Thus the produced states lie in the stipulated lower chamber.

For $b\in H$ the preceding universal threshold $b>64B_A^2$ is now replaced
by the factor-sensitive threshold
\[
 \boxed{b>8B_Aq_*.}
\]
The new theorem neither assumes the outside factor exists in every shell
nor asserts monotonicity of $D$ under an unspecified map.

\paragraph{Numerical certificate below the old sufficient threshold.}
Take $R=19$, $A=35$, and the quadratic residues
$H=\{1,4,5,6,7,9,11,16,17\}$. The actual anchor box covers $H$ and
$B_A=35$. At
\[
 p=558721=140\cdot3991-19,\qquad b=3991=13\cdot307,
\]
$b\in H$, $q_*=13$ and $T_H(b)=307>280$,
whereas $3991<78400=64B_A^2$. The cost at the selected residue is $35$,
with $(n,d)=(1,35)$. The theorem gives
\[
 (R,u,h,r,s,\lambda)=(19,15925,13,35,307,18)
\]
and ordered denominators
\[
 (139685,40137399198,4575924990).
\]
The prime is $121\pmod{840}$ and its primality is certified by complete trial
division. Three further examples and all minimizing pairs are in the JSON.
These are illustrations of a stronger criterion, not claims of previously
unknown ES coverage.

For a hypothetical false prime, every actual saturated anchor in every
shell has \emph{no} outside prime factor at all: without an angular restriction
one outside divisor and the congruence construction already yield a middle
state. For the stronger chamber-restricted counterfactual, the exact
necessary inequalities are instead
$t\le8c_A(-t^{-1})$ for every available outside $t$ (for the $1/8$
subchamber). The two counterfactuals must not be conflated.

\section{Transport through the marked geometric and lattice interfaces}
\subsection{Scale, Cayley coordinates, and all factor labels}
For a rational atlas state $q=(a,Y,Z)$ retain
\[
 A=(-p,4a,4Y,4Z),\qquad e_3(A)=0.
\]
The cleared state $(pD,Dq)$ has Cayley vector $DA$. Thus its projective
point, all four face root triples, and their factor labels agree with those
of the original rational candidate:
\[
 t_{ij}=-4A_i/A_j\quad(j\ne i),\qquad \sum_{j\ne i}t_{ij}=4.
\]
The affine parameter $p$ or the equivalent full scale must remain available.
$D$ can then be retained directly or recomputed from the rational denominators.
It is not recoverable from branch labels alone.
The face inverse is $A_j=-4A_i/t_{ij}$ with $A_i$ retained.

For a flag $(i,j,k,l)$ put $r=t_{ij},s=t_{ik},t=t_{il}$,
$d=-\tfrac14(r-s)(r-t)$ and $\eta=s-t$. The unnormalized factors
\[
 \ell=U-rV,\qquad Q=-\tfrac14(U-sV)(U-tV)
\]
are retained for every flag, with resultant $d$.
If $d\ne0$, the normalized pair is $(\ell/d,dQ)$ and the inherited Fable
source coordinates are
\[
 (d^{-1},-r-d,5d^2+3rd+\tfrac12d^3).
\]
The derivative, $\eta$, scale and flag are not deleted. If $d=0$,
the resultant-one chart is not invoked: the original factors and their
multiplicity labels remain. In particular the rejected $M,u=a$ states
are not thrown away to make the normalized chart look complete.
The selected computations carry every one of 24 flags, including the
zero-resultant strata, through the stated return.

The four-face resolvent from the preceding package is
$b_0=A_0A_1+A_2A_3$ and its two matching analogues. Clearing gives
$b_i(DA)=D^2b_i(A)$. This degree-two scale is retained. It does not identify
the three resolvent values with the three denominator values.

\subsection{A genuine failed rational candidate with integral geometric lifts}
For $p=241,a=64,u=1,E$, one has $R=15$, $D=3$ and
\[
 q=(64,3904/3,14701/3),\qquad
 X=(192,3904,14701),\qquad n=723.
\]
The integer $X$ solves ES at $723$, not at $241$. Nevertheless both records
have the same projective Cayley point and the same labelled face cubics.
The certificate supplies all normalized simple factor routes for this failed
candidate, not just routes for already successful states.

\subsection{Leech fibres and the full Albert correction}
Retain the preceding integral map $P:\Lambda\to\Z^3$ and its specified
section $\sigma(t)=\sum t_i b_i$, where the numerator of $b_0$ is
\[
(-3,1,-1,1,1,-1,-1,1,1,1,1,1,1,-1,1,1,1,1,-1,-1,-1,-1,1,1)
\]
and $b_1=Cb_0,b_2=C^2b_0$, each divided by $\sqrt8$.
The three octads, code membership and identities $P(b_i)=e_i$ are checked
exactly in the standalone package. The infinite section follows by
linearity and closure of the classical Leech lattice.

The clearing denominator is recoverable even from the integral trace, provided
the original prime is retained:
\[
 D=\frac{4X_0X_1X_2}{p\sum_{i<j}X_iX_j}.
\]
It then recovers $q=X/D$, $a$, $R$, and $u$ by the candidate inverse.
The channel is distinguished by $v_p(Y)=0$ in $E$ and $v_p(Y)>0$ in $M$.
Thus the map from canonical candidates to their ordered cleared triples,
and hence to their chosen Leech lifts, is injective for fixed $p$.
The complete scaled-candidate lift is
\[
 (\alpha,k)\longmapsto \lambda_\Lambda=\sigma(Dq)+k,
 \qquad k\in\ker P.
\]
Its inverse uses $q=P\lambda_\Lambda/D$, the retained $p,R,c$ and the
candidate inverse, then $k=\lambda_\Lambda-\sigma(Dq)$.
These are maps on the explicitly specified image, not on arbitrary
unmarked lattice vectors. The supplied tests use $k=0$ but retain the
full fibre theorem. The all-octad inverse is checked on these vectors too;
its nonzero-syndrome classes are not silently admitted.

The signed Golay--Wilson map in \texttt{lineage/aligned\_isometry.json}
identifies the full numerator vector with ordered octonions $z_0,z_1,z_2$.
Its explicit signed inverse is checked. The previously established
infinite-lattice image theorem is inherited, not newly audited here. The
checker verifies the signed dictionary against the pinned matrix on all
576 entries, tests the actual joint Wilson membership conditions on each
returned state, and computes the full octonion cubic using the retained
Fano multiplication convention \cite{wilson}.
For $X=P\lambda_\Lambda$, $n_i=N(z_i)$ and
$\nu=2\Re((z_0z_1)z_2)$, retain
\[
 \widehat X=\begin{pmatrix}X_0&z_2&\bar z_1\\\bar z_2&X_1&z_0\\
 z_1&\bar z_0&X_2\end{pmatrix}.
\]
Then the earlier exact identity becomes
\[
 4N_J(\widehat X)-pS_J(\widehat X)
 =2p\delta_\alpha\sum_{i<j}X_iX_j+\sum_i(p-4X_i)n_i+4\nu.
\]
The seven geometric test states include the numerical values of
$N_J,S_J,\nu$, and both sides of this identity. The diagonal denominator
defect and the off-diagonal cubic correction are separate terms. Neither an arbitrary lattice isometry nor a split-Tits
transformation is declared to preserve the marked compact cubic.

For any genuine bijection of retained states $f$, multiplication by
$\delta\circ f^{-1}$ is conjugate to multiplication by $\delta$ on the original
function algebra. Thus success kernels and negative inverse certificates
transport exactly. A many-to-one projection must retain its fibre, and a
map changing the arithmetic state must recompute $D$; functorial notation
alone does not prove an improvement.

\section{The infinite trace-space counterfactual and its analytic boundary}
Since $P^{-1}(t)=\sigma(t)+K$, $K=\ker P$, the positive trace space is
explicitly $\Z_{>0}^3\times K$. Multiplication by
$F_p(P\lambda)^2$ has a nonzero kernel exactly when ES holds at $p$.
Under the false-prime hypothesis all values are positive integers, so its
pointwise inverse $1/F_p(P\lambda)^2$ exists and is bounded by one.
Under the true-prime alternative every successful trace gives an entire
rank-21 lattice fibre of zeros. More precisely, on
$\ell^2(\{\lambda\in\Lambda:P\lambda\in\Z_{>0}^3\})$ with counting measure,
its domain is
\[
 \mathcal D(\mathcal H_p)=\{f:\sum_\lambda
    F_p(P\lambda)^4|f(\lambda)|^2<\infty\}.
\]
Coordinatewise multiplication by the nonnegative real values
$F_p(P\lambda)^2$ is self-adjoint on this maximal domain: the adjoint
condition tested against each coordinate vector forces the same multiplier
and the same square-summability requirement. In the false case its inverse
is the everywhere-defined multiplier $F_p(P\lambda)^{-2}$, with norm at
most one, and $\mathcal H_p\ge I$. In the true case its zero eigenspace is
infinite-dimensional, since an entire translate of the rank-21 kernel is
present. This is an exact operator encoding of the arithmetic alternative,
not an independently constructed differential operator or an asserted
analytic solution of the conjecture.

The same alternative is expressed by the previously marked theta series:
\[
 \mathcal Z_p(q)=
 \sum_{\substack{t\in\Z_{>0}^3\\F_p(t)=0}}
 q^{\|t\|^2/8}\Theta_{K+r(t)}(q),\qquad 0<q<1.
\]
Every summand is positive, so $\mathcal Z_p(q)=0$ exactly in the false
case. This identifies the analytic quantity to control; it does not
calculate it independently of arithmetic occupancy.

There is a useful external restriction on any genuine counterexample.
Bright--Loughran's Theorem 1.8 proves that the positive-real/integral-adic
Brauer--Manin set is nonempty for every parameter \cite{bl}. A false
prime would therefore have a nonempty such adelic set but no global positive
integer witness. This rules out one proposed emptiness explanation, not all
possible arithmetic obstructions. It also does not turn local compatibility
inside different exponent vectors into one simultaneous bounded lift.

\section{The two research directions as one program}
The positive track tries to construct a zero of $\delta_\alpha$. The negative
track assumes $\delta$ is a unit and derives its mandatory arithmetic,
concentration, geometric and integral-return properties. A contradiction
must violate one of those properties by a proved calculation.
The supplied positive tests are Theorem~\ref{thm:energy} and
Theorem~\ref{thm:cofactor}; the exact same-shell obstruction is
Proposition~\ref{prop:transport}.

A next useful theorem would prove that every full atlas has a shell whose
integer energy falls below its counterfactual minimum, or construct a
well-founded cross-shell move that strictly lowers $D$ until $D=1$.
No such universal decrease is assumed. The inverse certificate and the
restricted examples make a proposed move testable on failures as well as
successes. A proof of either universal assertion would finish the positive
prime cases; the present tranche proves neither assertion.

\section{Certificates, checks, and nonclaims}
The standard-library checker exhausts every full canonical candidate at
the following seven primes:
\begin{center}\begin{tabular}{r|rrrrrrr}
$p$&5&13&37&97&241&1201&2521\\
leaves&12&64&312&1216&3964&29600&73060
\end{tabular}\end{center}
For $1201$ the full counts are $8E+8M$, not the $7E+8M$ of the restricted
$R<p$ atlas; at $2521$ they are $7E+6M$, not $6E+6M$.
These are marked first-denominator counts, not sorted-triple counts.
All assertions about rational denominators, scaling, inverse coordinates
and repeated-root strata are checked on every candidate.

The two negative examples are labelled \emph{restricted}; the validator
rejects promoting them to full certificates. It also rejects a removed
leaf, an incorrect clearing denominator, and a fabricated full negative
certificate at $p=5$. The exact energy autocorrelations are recomputed
independently of their histogram squares. The largest-outside-divisor
formula is compared with exhaustive divisor enumeration for every unit
cofactor $b\le3000$ at $R=19$. All 24 factor labels are retained for seven
selected rational or integral states, including one repeated-root state.
No actual full negative certificate is found.

Run
\begin{verbatim}
python dual_es.py --self-test --out certificates
python -O dual_es.py --self-test --out certificates_optimized
\end{verbatim}
The optimized run has identical mathematical outputs. Use
\texttt{--prime p} to produce a witness or a fully verified negative
certificate, and \texttt{--verify-negative file.json} to check a proposed
full negative certificate. Trial division and complete enumeration are
finite but are not advertised as polynomial in $\log p$.
The source Giuga program was read, not fully executed; its expensive
power-sum and symbolic Cayley--Dickson suites are not part of these receipts.
No new numerical verification bound, independent reviewer, formal Lean
build, or proof of the global conjecture is claimed.

\begin{thebibliography}{9}
\bibitem{wilson} Wilson, R. A. (2009). Octonions and the Leech lattice.
\emph{Journal of Algebra, 322}, 2186--2190. Source PDF \S2--\S4:
Fano multiplication convention, joint lattice conditions, and coordinate comparison.
\url{https://maths.qmul.ac.uk/~raw/pubs_files/octoLeech1rev.pdf}.
\bibitem{yamamoto} Yamamoto, K. (1965). On the Diophantine equation
$4/n=1/x+1/y+1/z$. \emph{Memoirs of the Faculty of Science, Kyushu
University, Series A, Mathematics, 19}(1), 37--47.
Lemma 2, equation (4), p.~38.
\url{https://www.jstage.jst.go.jp/article/kyushumfs/19/1/19_1_37/_pdf}.
\bibitem{mv} Monks, M., \& Velingker, A. (2008, September 23).
\emph{On the Erd\H{o}s--Straus conjecture: Properties of solutions to its
underlying Diophantine equation}. Theorem 2.2 and Corollary 2.3, p.~4;
proof pp.~5--6. Only the reduced fraction case is used.
\url{https://mathematicalgemstones.com/maria/papers/ErdosStraus.pdf}.
\bibitem{et} Elsholtz, C., \& Tao, T. (2013). Counting the number of solutions
to the Erd\H{o}s--Straus equation on unit fractions.
\emph{Journal of the Australian Mathematical Society, 94}(1), 50--105.
\S2, Propositions 2.2 and 2.6. \url{https://arxiv.org/abs/1107.1010}.
\bibitem{kellner} Kellner, B. C. (2004). \emph{The equivalence of Giuga's
and Agoh's conjectures}. arXiv:math/0409259v1.
\S3, Lemma 3.1 and Theorem 3.5, pp.~4--5.
\url{https://arxiv.org/pdf/math/0409259}.
\bibitem{bl} Bright, M., \& Loughran, D. (2020). Brauer--Manin obstruction
for Erd\H{o}s--Straus surfaces. \emph{Bulletin of the London Mathematical
Society, 52}(4), 746--761. arXiv v2, Theorem 1.8, p.~4; proof \S3.8,
p.~14. \url{https://arxiv.org/abs/1908.02526}.
\bibitem{foundation} The Clankers. (2026). \emph{Exact bounded congruence
transport}, expanded workbench edition. Pinned repository commit
\texttt{9a3ddfe8b38a424f1ff0426201a1f4f6f9d2772f}.
\texttt{counterexample\_sieve/all\_shell\_no\_hit.tex},
Theorem \texttt{thm:shellwise-no-hit}, equations (1)--(10d);
\texttt{first\_two\_shell\_sieve.tex}, residual-three and residual-seven
theorems. Source-specific human credit remains attached to that edition.
\url{https://github.com/KokunoYumeto/erdos-straus-foundation}.
\end{thebibliography}
\end{document}
